\section*{अध्याय \ref{ch:b2:alpine-orogen} --- एक पर्वत-पट्टी की शारीरिकी: आल्प्स}

\begin{solution}{exo:b2:alpine-orogen:1}
(क) $e = 1 - 27.0/36.0 = 0.25$, यानी \qty{25}{\percent}; $S = 36.0 - 27.0 =
\qty{9.0}{km}$। (ख) परत ने अपनी लंबाई बनाए रखी हो (उसके अनुदिश न तनाव, न पतलापन),
काट परिवहन-दिशा के समांतर हो, और पिन रेखाएँ उन स्थानों को चिह्नित करें
जहाँ स्तर एक-दूसरे पर से नहीं सरके।
\end{solution}

\begin{solution}{exo:b2:alpine-orogen:2}
(क) एक \omterm{def:b2:alpine-orogen:klippe}{क्लिप}: नाइस अपने \omterm{def:b2:alpine-orogen:nappes}{स्वस्थानिक} आधार पर टिका किसी \omterm{def:b2:alpine-orogen:nappes}{नैप} का एक पृथक अवशेष है;
अपरदन ने उसके चारों ओर से चादर का शेष भाग हटा दिया है। (ख) एक \omterm{def:b2:alpine-orogen:klippe}{विवर्तनिक खिड़की}: अपरदन ने
नाइस \omterm{def:b2:alpine-orogen:nappes}{नैप} को काटकर उसके आधारी अधिक्षेप के नीचे की छोटी उम्र की चट्टानों को अनावृत कर दिया है, और इसलिए वह अधिक्षेप
हर ओर खिड़की से दूर की ओर नत है।
\end{solution}

\begin{solution}{exo:b2:alpine-orogen:3}
(क) अल्ट्रामैफ़िक संकुल (ऊपरी प्रावार); गैब्रोई संकुल (निचली महासागरीय भूपर्पटी);
चादरी डाइक-संकुल (मध्य भूपर्पटी); ज्वालामुखीय संकुल, तकिया-रूपी (ऊपरी भूपर्पटी)।
(ख) चादरी डाइक-संकुल। (ग) चौड़े महासागरों की अधिकांश भूपर्पटी क्षेपण में खप जाती है;
जो \omterm{def:b2:alpine-orogen:ophiolite}{ओफ़ियोलाइट} स्थल पर बचे रहते हैं वे अधिकतर छोटी महासागरीय द्रोणियों में या द्वीप-चापों के
बीच बने, और टेथिस वाले वहाँ जहाँ कोई महाद्वीप तब तक पतला हुआ जब तक एक छोटा महासागर नहीं खुल गया।
\end{solution}

\begin{solution}{exo:b2:alpine-orogen:4}
(क) मोंते रोसा और ग्रां सें बर्नार \omterm{def:b2:alpine-orogen:nappes}{नैप}: यूरोपीय पक्ष (ब्रियांसोनै और
यूरोपीय निष्क्रिय किनारा); त्सेरमाट--सास और कॉम्बैं मंडल: पिएमोंट--लिगुरिया महासागर; दाँ
ब्लांश तंत्र: अपूलिया का किनारा। (ख) अपूलिया, जो शीर्ष पर समाप्त होती है। (ग) ऊपरी प्लेट का
क्षेपण नहीं होता, और निचली प्लेट से खुरची गई पतली परतें उसके नीचे उसी क्रम में ढेर होती हैं जिसमें
निचली प्लेट पहुँची: दूर का महासागर पहले, निकट का किनारा अंत में और सबसे नीचे। ढेर का बाद का वलन,
या क्रम-बाह्य अधिक्षेप, इस क्रम को बिगाड़ सकते हैं।
\end{solution}

\begin{solution}{exo:b2:alpine-orogen:5}
(क) $2\pi G\,\Delta\rho = 2\pi\times6.674\times10^{-11}\times500 =
\qty{2.097e-7}{s^{-2}}$; $h = 1.70\times10^{-3}/2.097\times10^{-7} = \qty{8.1}{km}$।
(ख) $2.097\times10^{-7}\times25\,000 = \qty{5.24e-3}{m.s^{-2}}$, यानी लगभग
\qty{-524}{mGal}। (ग) जड़ कोई अनंत पट्टी नहीं है: वह कुछ दसियों किलोमीटर
गहरा और केवल पट्टी जितना चौड़ा एक पिंड है, इसलिए वह सतह के किसी बिंदु को उतनी ही मोटाई की
अनंत तक फैली पट्टी से कम खींचता है; इव्रिया मंडल जैसे सघन पिंड
धनात्मक विसंगतियाँ जोड़ते हैं, और गहराई पर घनत्व-अंतर मान्य मान से कम हो सकता है।
\end{solution}

\begin{solution}{exo:b2:alpine-orogen:6}
(क) $164 - 44 = \qty{120}{Myr}$। (ख) $z = P/(\bar\rho g)$:
$2.8\times10^9/(3000\times9.8) = \qty{95.2}{km}$; $2.8\times10^9/(3300\times9.8) =
\qty{86.6}{km}$। (ग) ``\qty{90}{km} से अधिक'' इसी परास के भीतर है: गहराई अधिभार के लिए चुने गए
माध्य घनत्व पर उतनी ही निर्भर करती है जितनी दाब पर।
\end{solution}

\begin{solution}{exo:b2:alpine-orogen:7}
(क) लंबाई $L$ का एक पार्श्व, जिसकी नति $\theta$ है, क्षैतिज रूप से $L\cos\theta$ में फैलता है; सभी पार्श्वों
की नति एक ही है, अतः $l = l_0\cos\theta$ और $e = 1 - l/l_0 = 1 - \cos\theta$।
(ख) $\cos\theta = 0.70$: $\theta = \qty{45.6}{\degree}$। (ग) $1 - \cos\qty{40}{\degree} =
\qty{23.4}{\percent}$, जैसा शेवरॉन चित्र में; $1 - \cos\qty{60}{\degree}
= \qty{50}{\percent}$।
\end{solution}

\begin{solution}{exo:b2:alpine-orogen:8}
(क) $l_0 = 118/2.0 = \qty{59.0}{km}$; $S = 59.0 - 42 = \qty{17.0}{km}$; $e = 17.0/59.0 =
\qty{28.8}{\percent}$। (ख) $A_e = 118 - 42\times2.0 = \qty{34.0}{km^2}$; $d = 34.0/17.0 =
\qty{2.0}{km}$। (ग) हाँ: \omterm{def:b2:alpine-orogen:decollement}{डेकोलमां} \qty{2.0}{km} मोटे आवरण के आधार पर, उसके शीर्ष के नीचे,
स्थित है; आवरण अपने आधार पर अलग हुआ है।
\end{solution}

\begin{solution}{exo:b2:alpine-orogen:9}
(क) एयरी जड़ $4.5h$ है: \qty{46.5}{\degree} N पर $4.5\times2.74 = \qty{12.3}{km}$, यानी
\qty{42.3}{km} पर मोहो; \qty{45.55}{\degree} N पर $4.5\times0.65 = \qty{2.9}{km}$, यानी
\qty{32.9}{km} पर मोहो। (ख) ऊँची भूमि के नीचे मोहो एयरी के पूर्वानुमान से \qty{6.0}{km} उथला है;
पो मैदान के किनारे के नीचे वह \qty{28.3}{km} अधिक गहरा है। (ग) जड़ की
प्रतिपूर्ति स्थानीय नहीं है: वह दक्षिण की ओर विस्थापित है और उच्चावच की आवश्यकता से बड़ी है।
प्लेट अपने ऊपर के भारों के नीचे समग्र रूप से झुकती है (तृतीय वर्ष के एक अध्याय का क्षेत्रीय
आलंब), और उन भारों में से एक, इव्रिया पिंड, स्थलाकृति है ही नहीं।
\end{solution}

\begin{solution}{exo:b2:alpine-orogen:10}
(क) $5.1/11.3 = 0.45$: आधे से कम। (ख) $5.1 + 6.2 = \qty{11.3}{km}$। जुरैसिक का शीर्ष
ऊपरी विलगन-तल के ऊपर स्थित है; उस विलगन-तल के अनुदिश ले जाया गया \qty{6.2}{km}
काट के वलनों के नीचे से गुज़रा और आगे उत्तर की संरचनाओं में खप गया, इसलिए
ऊपरी परत काट के भीतर उतनी कभी छोटी नहीं हुई। (ग) एक अकेली चिह्नक परत केवल उन्हीं
स्तरों का लघुकरण अंकित करती है जिनसे वह जुड़ी है; उसके ऊपर या नीचे का हर विलगन-तल विस्थापन को काट में
ला या उससे बाहर ले जा सकता है। कई परतों का संतुलन कीजिए, और क्षेत्रफलों से जाँचिए।
\end{solution}

\begin{solution}{exo:b2:alpine-orogen:11}
(क) \qty{100}{km} के विस्थापन के लिए $\sqrt{200^2 + 100^2} = \qty{223.6}{km}$;
\qty{50}{km} के लिए $\sqrt{200^2 + 50^2} = \qty{206.2}{km}$; \qty{150}{km} के लिए
$\sqrt{200^2 + 150^2} = \qty{250.0}{km}$। (ख) उत्तर--दक्षिण से $\arctan(100/200) = \qty{26.6}{\degree}$,
एक अभिसरण-रेखा जो रोज़नबर्ग और सहकर्मियों द्वारा उद्धृत N\qty{150}{\degree}E
के निकट है। (ग) विस्थापन में $\pm\qty{50}{km}$ की अनिश्चितता है, जो कुल
अभिसरण को किसी भी ओर लगभग \qty{20}{km} खिसका देती है; परिणाम का तीसरा अंक
उस अनिश्चितता से बहुत नीचे है।
\end{solution}

\begin{solution}{exo:b2:alpine-orogen:12}
(क) $210/12.9 = \qty{16.3}{Myr}$ से $210/8.5 = \qty{24.7}{Myr}$ तक (\unit{km} को
\unit{mm.yr^{-1}} से भाग देने पर मिलियन वर्ष मिलते हैं)। (ख) $4\times60 = \qty{240}{km}$ से अधिक;
द्रोणी का कम से कम $240 - 60 = \qty{180}{km}$ अब ऐल्पाइन अधिक्षेपों के नीचे है या अपरदित हो
चुका है। (ग) द्रोणी अधिक्षेप-कील के सामने बनती है और उसके साथ स्थानांतरित होती है: उसका भीतरी
भाग सदा उसे भार देने वाले अधिक्षेपों द्वारा अधिक्षेपित होता रहता है, और उसका सबसे पुराना भराव
अधिक्षेपों में फँस जाता है, जैसा उप-ऐल्पाइन \omterm{def:b2:alpine-orogen:foreland}{मोलास} है।
\end{solution}

\begin{solution}{pb:b2:alpine-orogen:1}
\noindent\textbf{भाग I।}

\textbf{1.} $3.0/\cos\qty{35}{\degree} = \qty{3.66}{km}$; $2.5/\cos\qty{50}{\degree} =
\qty{3.89}{km}$; $4.0/\cos\qty{40}{\degree} = \qty{5.22}{km}$।

\textbf{2.} $8.00 + 2\times3.66 + 5.00 + 2\times3.89 + 6.0/\cos\qty{10}{\degree} +
2\times5.22 + 10.0/\cos\qty{5}{\degree} = 8.00 + 7.32 + 5.00 + 7.78 + 6.09 + 10.44 + 10.04 =
\qty{54.67}{km}$; दोहराए गए हिस्से सहित, $l_0 = \qty{60.7}{km}$।

\textbf{3.} अधिक्षेप-चादर के नीचे की परत वही परत है, जिसे अधिक्षेप चादर के नीचे ले गया:
अधिक्षेपण से पहले वह शेष के साथ एक सीध में थी। उसे छोड़ देने से परत
बहुत छोटी पुनर्स्थापित होती।

\textbf{4.} $e = 1 - 48.0/60.7 = 0.209$, यानी \qty{20.9}{\percent}; $S = 60.7 - 48.0 =
\qty{12.7}{km}$।

\textbf{5.} $60.7 - 60.2 = \qty{0.5}{km}$, $0.5/60.7 = \qty{0.8}{\percent}$: मापन की
परिशुद्धता के भीतर; काट स्वीकार्य है।

\textbf{6.} $60.7 - 54.0 = \qty{6.7}{km}$, \qty{11}{\percent}: एक ही चादर की दो परतें
इतनी भिन्न लंबाइयों तक पुनर्स्थापित नहीं हो सकतीं। काट जैसी खींची गई है वैसी स्वीकार्य नहीं होती: कोई
संरचना छूटी या ग़लत खींची गई, कोई पिन रेखा ग़लत स्थान पर, काट के तल से बाहर गति,
या दोनों परतों के बीच कोई विलगन-तल।

\noindent\textbf{भाग II।}

\textbf{7.} $146/2.4 = \qty{60.8}{km}$।

\textbf{8.} यह \qty{60.7}{km} से \qty{0.1}{km} भिन्न है: रेखा-लंबाई और क्षेत्रफल सहमत हैं।

\textbf{9.} $A_e = 146 - 48.0\times2.4 = 146 - 115.2 = \qty{30.8}{km^2}$।

\textbf{10.} $d = 30.8/12.7 = \qty{2.43}{km}$, जो आँकड़ों की परिशुद्धता तक आवरण के
\qty{2.4}{km} के बराबर है: \omterm{def:b2:alpine-orogen:decollement}{डेकोलमां} आवरण का आधार है।

\textbf{11.} $160/2.4 = \qty{66.7}{km}$, रेखा-लंबाई से \qty{6.0}{km} अधिक। आवरण के आधार पर
वाष्पनिक निक्षेपों का मोटा होना, जो अपनतियों के क्रोडों में बहते हैं, लंबाई जोड़े बिना क्षेत्रफल
जोड़ता है; तब स्थिर मोटाई वाला क्षेत्रफल संतुलन विफल हो जाता है।

\textbf{12.} आवरण अपने आधार पर एक \omterm{def:b2:alpine-orogen:decollement}{डेकोलमां} के ऊपर चला और आधारशैल
सम्मिलित नहीं था। पट्टी के नीचे आधारशैल का शीर्ष समतल या हल्का और एकसमान नत होना चाहिए,
अपने ऊपर के अधिक्षेपों से अटूट।

\noindent\textbf{भाग III।}

\textbf{13.} $22 + 18 + 25 = \qty{65}{km}$; $S = 65 - 15 = \qty{50}{km}$।

\textbf{14.} $e = 50/65 = \qty{76.9}{\percent}$, जबकि अग्रभूमि पट्टी में
\qty{20.9}{\percent}: \omterm{def:b2:alpine-orogen:nappes}{नैप} कहीं अधिक छोटे हुए हैं।

\textbf{15.} $30 - 22 = \qty{8}{km}$ ($e = \qty{26.7}{\percent}$)। स्वयं आधारशैल
छोटा हुआ है: \omterm{def:b2:alpine-orogen:skin}{मोटी-त्वचा विवर्तनिकी}।

\textbf{16.} अग्रभूमि पट्टी $12.7/(11 - 4) = \qty{1.81}{km/Myr}$, यानी
\qty{1.81}{mm} प्रति वर्ष; \omterm{def:b2:alpine-orogen:nappes}{नैप} $50/(30 - 20) = \qty{5.0}{mm}$ प्रति वर्ष।

\textbf{17.} \omterm{def:b2:alpine-orogen:nappes}{नैप} (30 से \qty{20}{Ma}) अग्रभूमि पट्टी (11 से
\qty{4}{Ma}) से पहले विरूपित हुए: विरूपण बाहर की ओर, अग्रभूमि की दिशा में, स्थानांतरित हुआ। यही
आल्प्स का क्रम है: ओलिगोसीन और मायोसीन में ग्लारुस अधिक्षेप, मायोसीन में बाह्य पुंजक,
उत्तर मायोसीन और प्लायोसीन में जूरा।

\textbf{18.} $12.7 + 50 + 8 = \qty{70.7}{km}$।

\noindent\textbf{भाग IV।}

\textbf{19.} $\qty{0.5}{cm.yr^{-1}}\times\qty{40}{Myr} = 0.005\times4\times10^7 =
\qty{2.0e5}{m} = \qty{200}{km}$।

\textbf{20.} $200 + 18 = \qty{218}{km}$।

\textbf{21.} केवल वही गिना जाता है जो संरक्षित है या अनुमित है: महासागर के कुछ भागों का
\qty{85}{Ma} से पहले या माध्य दर से तेज़ क्षेपण हुआ हो सकता है, ओफ़ियोलाइट-पट्टियाँ
अपरदित हुई या काट से बाहर ले जाई गई हो सकती हैं, और काट से तिर्यक अभिसरण
ऐसा महासागर-तल लाता है जिसे काट नहीं देखती।

\textbf{22.} आंतरिक मंडल $60 - 20 = \qty{40}{km}$ ($e = \qty{66.7}{\percent}$); प्रति-कील
$70 - 48 = \qty{22}{km}$ ($e = \qty{31.4}{\percent}$)।

\textbf{23.} $12.7 + 50 + 8 + 40 + 22 = \qty{132.7}{km}$; यह मध्य आल्प्स के अनुमानों के परास के भीतर है,
यानी लगभग 125 से \qty{160}{km} टक्कर-जनित लघुकरण (स्ट्राइक-स्लिप सरकन जोड़ने से
पहले)।

\textbf{24.} पुनर्स्थापित चौड़ाई $153 + 132.7 = \qty{285.7}{km}$; पुनर्स्थापित भूपर्पटी
$285.7\times30 = \qty{8571}{km^2}$; वर्तमान भूपर्पटी $153\times45 = \qty{6885}{km^2}$। ग़ायब:
$8571 - 6885 = \qty{1686}{km^2}$, पुनर्स्थापित भूपर्पटी का \qty{19.7}{\percent}; यह
अपरदित हो गई है (और \omterm{def:b2:alpine-orogen:foreland}{अग्रभूमि द्रोणियों} में बहा दी गई है) या निचली भूपर्पटी के साथ प्रावार में
क्षेपित हो गई है। इन हानियों के बिना भूपर्पटी औसतन $8571/153 = \qty{56}{km}$ मोटी होती।
पूरी काट पर कुल टक्कर-जनित लघुकरण
\textbf{\qty[text-series-to-math=true]{133}{km}} है।
\end{solution}
